% Nouveau EGA V (cote 25): archive cote 25, site page 160 (PDF page 161).
% English translation of the checked French annotated master.
% AI-translated by GPT-6 Astra (OpenAI, ChatGPT Pro mode).
\scanpage{25/160}
\mspage{\hand{3}}

proposition seems correct, and should be proved in an entirely analogous way
by showing that for a given $x\hand{\in} X$\var{\inserted[typed]{$^{(k)}$}}{$^{(k)}$}{insertion}{ega25_25_160-v0001}, there is a hypersurface \var{\inserted[typed]{$H$}}{$H$}{insertion}{ega25_25_160-v0002} of given degree $n\hand{\geqslant} 2$
which is tangent to X at x, and such that x is an elementary quadratic
point of $X$\hand{$\cap$}$H$, of dimension one less .... Be quite blunt about
it, by pointing out that if k \uline{has} char.2 and X is
connected \var{\inserted[typed]{nonempty}}{nonempty}{insertion}{ega25_25_160-v0003} of odd dimension, then the conditions of 8.8. \uline{are not}
satisfied, i.e. L/K is necessarily inseparable (of degree
2 exactly, I believe, if $X\hand{\subset} P^r$, but with no guarantee ...).

15.1. Let S be a prescheme, $\underline{E}$ a loc free Module of finite type
over S, $\hand{\phi}$ a section of $\underline{\mathrm{Sym\var{\xout{m}}{}{typewriter strike-out}{ega25_25_160-v0004}}}^n(\underline{E})$, i.e. an ``n-form on $\overset{\hand{\scriptscriptstyle\vee}}{\underline{E}}$'', where $n$\var{\over{$>$}{\unsure{$\geqslant$}}}{\unsure{$\geqslant$}}{writing over}{ega25_25_160-v0005}$0$.
Giving $\hand{\phi}$ is equivalent to giving a section of $\underline{O}_P(n)$ on $P=P(\underline{E})$
(ref to III), and thus defines a subprescheme $V(\hand{\phi})$ of P. We shall say that
the form $\hand{\phi}$ is \uline{smooth} if $V(\hand{\phi})$ is smooth and if, in addition, for every $s\hand{\in} S$,
$V(\hand{\phi})_s\hand{\neq} P_s$ (i.e. $\hand{\phi}(s)\hand{\in} \mathrm{Sym}^n_{k(s)}E\otimes k(s)$ is \var{\over{\unsure{$=$}}{$\neq$}}{$\neq$}{writing over}{ega25_25_160-v0006}\hand{\unsure{(}}$0$ ). One sees at once
that $\hand{\phi}$ is smooth iff for every $\unsure{x|s}$ $\hand{\in}$S, $\hand{\phi}(s)$ is smooth, \var{\xout{\unsure{and that over}}}{}{typewriter strike-out}{ega25_25_160-v0007}
\var{\xout{a field}}{}{typewriter strike-out}{ega25_25_160-v0008}(which reduces us to the case where S is the spectrum of a field), and
that the notion of a smooth form is invariant under change of the base
field (which reduces us, in particular, to the case of an alg
closed base field); these two properties can be summarized by saying that if $S'\hand{\to} S$ is
a surjective morphism, then $\hand{\phi}$ is smooth iff its inverse image $\hand{\phi}'$ over S'
is smooth. Of course, if E is a projective module of finite type over a
ring A, and $\hand{\phi}\hand{\in}\mathrm{Sym}^n_A(E)$, one says that $\hand{\phi}$ is smooth if .... When $E=A^{\unsure{N|n}+1}$,
giving $\hand{\phi}$ is equivalent to giving a \var{\xout{form}}{}{typewriter strike-out}{ega25_25_160-v0009} homogeneous polynomial of degree n
in variables $X_0,..,X_r$, and the Jacobian criterion immediately implies
that \var{\inserted{\unsure{if} $n>0$\unsure{?}}}{\unsure{if} $n>0$\unsure{?}}{insertion}{ega25_25_160-v0010} $\hand{\phi}$ is smooth iff the ideal generated by $\hand{\phi}$ and the $\hand{\dfrac{\partial\hand{\phi}}{\partial\textcolor{black}{X_i}}}$ contains
a power of the augmentation ideal $(X_0,\dots,X_r)$, i.e. contains a

\mspage{\hand{160}} \var{\margin{For $n=0$: one finds the invertible sections of $\underline{O}_S$ !}}{\placedmargin{For $n=0$: one finds the invertible sections of $\underline{O}_S$ !}}{marginal note placed}{ega25_25_160-v0011}
\tnote{25-160-1}{Be quite blunt about it renders the colloquial mettre les pieds dans le plat. The author's parenthetical I believe and with no guarantee are retained. Loc, alg and ref mean locally, algebraically and reference; sss is rendered iff. The unfinished smooth if clause and the questionable inserted if with its question mark remain.}
\tnote{25-160-2}{The canceled e forme contains a French agreement letter followed by form; the English variant retains the meaningful word form. The mathematical operator with its canceled m, the uncertain comparison at the degree bound, the doubtful parenthesis, the x/s choice and the N/n exponent all remain exact. The later variable list ends in r, as in the French.}
